\section{Referee-round positive closure: density-state nonlinear semigroup}
\label{sec:rr2-b4}

\subsection{Weighted density state and cylinder core}

Fix \(m>4\) and let
\[
 \mathcal X_m=\left\{
 f\in\mathcal P(\mathbb T^3\times\mathbb R^3):
 \int(1+|v|^m)\,df<\infty
 \right\}
\]
with the weighted bounded-Lipschitz metric.  Let \(\mathscr D\) be the algebra
of smooth cylinders
\[
 F(f)=\Phi(\langle\phi_1,f\rangle,\ldots,\langle\phi_k,f\rangle),
\]
where \(\Phi\in C_b^3\) and the \(\phi_j\) have weighted \(C^2\) bounds.

For the finite hard-sphere process define the nonlinear generator on the
empirical state by
\[
 H_\epsilon F
 =
 \mu_\epsilon^{-1}e^{-\mu_\epsilon F}
 L_\epsilon e^{\mu_\epsilon F}.
\]

\begin{theorem}[Core generator convergence]
\label{thm:rr2-b4-generator}
For every \(F\in\mathscr D\),
\[
 H_\epsilon F(f^\epsilon)\longrightarrow
 H F(f)
 =
 \mathbb H\left(f,\frac{\delta F}{\delta f},0\right)
\]
locally uniformly on moment balls, under the source-dependent prepared law.
The same conclusion holds with a bounded collision-work source.
\end{theorem}

\begin{proof}
Apply the deterministic Liouville generator to the exponential cylinder.
Free flights give the transport term.  At a binary contact the jump in the
exponent is \(\Delta(\delta F/\delta f)\), and summing contact fluxes gives
the exponential collision term.  Multiple simultaneous collisions form a
codimension-two set and have zero prepared measure.  The source-uniform
marked cluster theorem of B2 controls recollision errors and makes the
convergence uniform on moment balls.  B1 supplies the source-dependent
constraint saddle; no fixed zero-source replacement is used.
\end{proof}

\begin{lemma}[Exponential compact containment]
\label{lem:rr2-b4-containment}
For every \(T,M\) there is a compact \(K_{T,M}\subset\mathcal X_m\) such
that the empirical density exits \(K_{T,M}\) before time \(T\) with
probability at most \(e^{-M\mu_\epsilon}\).
\end{lemma}

\begin{proof}
Energy conservation gives the second moment.  Exponential velocity moments
of the prepared law and the short-time marked cluster expansion control the
higher weighted moment.  The exact balance controls time increments against
the countable test family defining the metric.  Exponential Chebyshev and
weighted Arzel\`a--Ascoli give the compact set.
\end{proof}

\subsection{Comparison for the functional Hamilton--Jacobi equation}

Consider
\[
 \partial_tU+HU=0,\qquad U(T,\cdot)=\Phi.
\]

\begin{theorem}[Comparison on weighted moment balls]
\label{thm:rr2-b4-comparison}
For bounded uniformly continuous terminal data and the Hamiltonian above,
bounded viscosity sub- and supersolutions satisfy comparison on every
invariant weighted moment ball.  Hence the equation defines a unique
nonlinear semigroup \(S_{s,t}\).
\end{theorem}

\begin{proof}
Use the weighted bounded-Lipschitz distance and the doubling functional
\[
 U(t,f)-V(t,g)-\alpha d_m(f,g)^2-\eta
 \big(\mathcal M_m(f)+\mathcal M_m(g)\big).
\]
The transport part is locally Lipschitz in the dual weighted norm.  The
collision exponential is convex in the cotangent variable and locally
Lipschitz on bounded cotangent cylinders; the moment penalty controls its
velocity tails.  The standard doubling argument gives a contradiction to a
positive maximum after first \(\alpha\to\infty\) and then \(\eta\downarrow0\).
Moment invariance removes boundary escape.
\end{proof}

\begin{theorem}[Nonlinear semigroup convergence]
\label{thm:rr2-b4-semigroup}
For \(\Phi\) in the bounded uniformly continuous class,
\[
 \frac1{\mu_\epsilon}\log
 \mathbb E_{\rm mc,a}
 \exp\{\mu_\epsilon\Phi(f_t^\epsilon)\}
 \longrightarrow
 S_{0,t}\Phi(f_0)
\]
locally uniformly on admissible initial moment balls.  The same limit is
obtained from the exact complete-history tower; thus history is lost in the
kinetic limit on this observable class.
\end{theorem}

\begin{proof}
Theorem~\ref{thm:rr2-b4-generator} gives convergence of upper and lower
nonlinear generators on the separating core.  Lemma~\ref{lem:rr2-b4-containment}
gives exponential compact containment.  Theorem~\ref{thm:rr2-b4-comparison}
gives uniqueness of the limiting Hamilton--Jacobi problem.  The nonlinear
semigroup convergence theorem then applies.  The complete-history and
density evaluations have the same finite-dimensional cylinder generators;
B2's propagation of chaos under every compact source ball makes their
difference exponentially negligible, proving loss of history.
\end{proof}

\subsection{Prepared diagonal limit}

\begin{theorem}[Microcanonical preparation before scaling]
\label{thm:rr2-b4-diagonal}
Let shell widths tend to zero slower than
\(\mu_\epsilon^{-1/2}\).  Then the microcanonical nonlinear semigroup is
governed by the constrained Hamiltonian
\[
 H_aF
 =
 \inf_\lambda
 \left\{
 H_{\rm gc}^{\lambda}F-\langle\lambda,a\rangle
 \right\}
 -
 \inf_\lambda
 \left\{
 H_{\rm gc}^{\lambda}0-\langle\lambda,a\rangle
 \right\},
\]
with the saddle reoptimized at every source.  The preparation limit and the
Boltzmann--Grad limit commute on the cylinder core.
\end{theorem}

\begin{proof}
B1 gives a uniform \(o(\mu_\epsilon)\) coefficient-extraction remainder for
the source-dependent saddle.  Divide by \(\mu_\epsilon\) and combine with
the locally uniform generator convergence.  Compactness of the saddle set
and strict constraint convexity make the infimum stable.  A diagonal
argument on a countable core yields the result.
\end{proof}

\subsection{Mechanical ray and diffusive tangent}

Let \(G\) be an additive mechanical work direction and let the preparation
field be \(\theta G\).  On the same cotangent ray define
\[
 \mathcal E_G^\theta(F)=
 \theta^{-1}\big(Q(\theta G+\theta F)-Q(\theta G)\big).
\]

\begin{theorem}[Ray-calibrated theta semigroup]
\label{thm:rr2-b4-theta}
If \(F\) is a work perturbation in the same cotangent work bundle as \(G\),
unit covariance and preparation consistency force the coefficient in the
certainty equivalent to be the same \(\theta\).  Under the Gaussian scaling
of B3, \(S_{s,t}\) converges to the risk-sensitive diffusion semigroup with
generator
\[
 \mathcal L u+\frac{\theta}{2}
 \langle Du,D\,Du\rangle.
\]
\end{theorem}

\begin{proof}
The calibration statement is the infinite-dimensional analogue of
the calibrated-work theorem of Paper A1: a simultaneous rescaling of the work
bundle rescales both the preparation covector and every perturbation, so the
dimensionless coefficient is unique.  B3 gives convergence of the
fluctuation field and of pressure Hessians.  Taylor expansion of the full
Hamiltonian on the \(\mu_\epsilon^{-1/2}\) scale leaves the quadratic
covariance term, while the uniform third-derivative bound makes the
remainder vanish.  Comparison identifies the limiting diffusion semigroup.
\end{proof}
